\section{Experimental Results}
\label{sec:results}

This section reports the final cross-basis benchmark collected in a single
run using the same machine, compiler configuration, Goldilocks backend, and
Criterion methodology for every representation.

The final comparison includes:

\begin{enumerate}
    \item the optimized Boolean-kernel evaluator;
    \item the blocked parallel monomial evaluator with block size \(4096\);
    \item serial Horner evaluation;
    \item parallel roots-of-unity Lagrange evaluation using batch inversion;
    \item serial roots-of-unity Lagrange evaluation using batch inversion.
\end{enumerate}

For every input size, all five methods represent the same deterministic
polynomial and are checked for equality at the evaluation point before
timing.

\subsection{Cross-basis scaling}

\cref{tab:final-basis-results} reports the Criterion central estimates for
all six tested sizes.

\begin{table}[ht]
\centering
\small
\begin{tabular}{rrrrrr}
\toprule
\(N\) &
Kernel &
Monomial par. &
Horner &
Lagrange par. &
Lagrange ser. \\
\midrule
\(2^{10}\) &
\(4.343\,\mu\mathrm{s}\) &
\(8.613\,\mu\mathrm{s}\) &
\(8.708\,\mu\mathrm{s}\) &
\(109.710\,\mu\mathrm{s}\) &
\(21.612\,\mu\mathrm{s}\) \\

\(2^{12}\) &
\(16.975\,\mu\mathrm{s}\) &
\(34.599\,\mu\mathrm{s}\) &
\(35.491\,\mu\mathrm{s}\) &
\(250.220\,\mu\mathrm{s}\) &
\(84.269\,\mu\mathrm{s}\) \\

\(2^{14}\) &
\(59.153\,\mu\mathrm{s}\) &
\(69.174\,\mu\mathrm{s}\) &
\(146.720\,\mu\mathrm{s}\) &
\(439.930\,\mu\mathrm{s}\) &
\(334.240\,\mu\mathrm{s}\) \\

\(2^{16}\) &
\(91.116\,\mu\mathrm{s}\) &
\(139.480\,\mu\mathrm{s}\) &
\(558.910\,\mu\mathrm{s}\) &
\(1.2503\,\mathrm{ms}\) &
\(1.3391\,\mathrm{ms}\) \\

\(2^{18}\) &
\(263.230\,\mu\mathrm{s}\) &
\(380.910\,\mu\mathrm{s}\) &
\(2.2440\,\mathrm{ms}\) &
\(4.5612\,\mathrm{ms}\) &
\(5.3679\,\mathrm{ms}\) \\

\(2^{20}\) &
\(1.0560\,\mathrm{ms}\) &
\(1.4116\,\mathrm{ms}\) &
\(9.0402\,\mathrm{ms}\) &
\(17.1990\,\mathrm{ms}\) &
\(21.3890\,\mathrm{ms}\) \\
\bottomrule
\end{tabular}
\caption{Final same-run arbitrary-point evaluation benchmark over the
Goldilocks field. Lower is better.}
\label{tab:final-basis-results}
\end{table}

The Boolean-kernel evaluator is the fastest method at every tested size in
this final sweep.

\subsection{Headline result at \(N=2^{20}\)}

For
\[
N=2^{20}=1{,}048{,}576,
\]
the optimized Boolean-kernel evaluator has Criterion interval
\[
[1.0009,\ 1.1132]\ \mathrm{ms}
\]
with central estimate
\[
\boxed{1.0560\ \mathrm{ms}}.
\]

The corresponding throughput is approximately
\[
\boxed{9.93\times10^8}
\]
input coefficients per second.

The parallel monomial baseline has central estimate
\[
1.4116\ \mathrm{ms}.
\]
Thus, in the final same-run comparison,
\[
\frac{1.4116}{1.0560}
=
\boxed{1.337\times}.
\]

Equivalently, the Boolean-kernel evaluator uses approximately
\[
1-\frac{1.0560}{1.4116}
\approx
\boxed{25.2\%}
\]
less wall-clock time than the blocked parallel monomial evaluator.

Against serial Horner evaluation,
\[
\frac{9.0402}{1.0560}
=
\boxed{8.56\times}.
\]

\subsection{Comparison with the Lagrange representation}

At
\[
N=2^{20},
\]
parallel arbitrary-point evaluation from the roots-of-unity Lagrange
representation requires
\[
17.199\ \mathrm{ms},
\]
while the serial version requires
\[
21.389\ \mathrm{ms}.
\]

Relative to the optimized Boolean-kernel evaluator, these correspond to
\[
\frac{17.199}{1.0560}
=
\boxed{16.29\times}
\]
and
\[
\frac{21.389}{1.0560}
=
\boxed{20.25\times},
\]
respectively.

This result concerns \emph{arbitrary-point evaluation from the stored
representation}.  It should not be interpreted as evidence that
roots-of-unity evaluation form is generally inferior in a proof system.
Evaluation form is highly advantageous for other workloads, including
pointwise operations and FFT/NTT-based encoding.  The comparison instead
shows that an arbitrary off-domain evaluation has a very different cost
profile in the three representations considered here.

\subsection{Scaling relative to the kernel evaluator}

The ratio between each competitor and the Boolean-kernel evaluator changes
with \(N\).  For the fair parallel monomial baseline, the measured ratios
are approximately

\[
\begin{array}{c|cccccc}
N &
2^{10} &
2^{12} &
2^{14} &
2^{16} &
2^{18} &
2^{20}
\\
\hline
T_{\mathrm{mono,par}}/T_{\mathrm{kernel}}
&
1.98 &
2.04 &
1.17 &
1.53 &
1.45 &
1.34
\end{array}
\]

The non-monotonic behavior reflects the interaction between subtree
parallelism, Rayon scheduling, cache behavior, and the fixed block size.
It is therefore more appropriate to report the measured curve than to
infer a constant hardware-independent speedup factor.

For serial Horner evaluation, the ratio grows much more strongly with
problem size:
\[
2.00,\ 2.09,\ 2.48,\ 6.13,\ 8.52,\ 8.56.
\]
This is consistent with the fact that the kernel tree exposes substantial
intra-polynomial parallelism while Horner evaluation retains a long
sequential dependency chain.

\subsection{Best observed optimization run versus final comparison}

During optimization of the Boolean-kernel evaluator alone, the best
observed \(N=2^{20}\) run was
\[
\boxed{892.57\ \mu\mathrm{s}},
\]
with interval
\[
[884.25,\ 902.65]\ \mu\mathrm{s}.
\]

The final multi-method sweep measured the same optimized kernel evaluator at
\[
1.0560\ \mathrm{ms}.
\]

We deliberately use the latter number for the cross-basis headline ratio,
because it was collected in the same benchmark campaign as all competing
representations.  The lower \(892.57\,\mu\mathrm{s}\) measurement is retained
as an optimization result and evidence of run-to-run variability, but it is
not mixed with competitor timings from a different run to construct a
speedup claim.

This distinction prevents accidental cherry-picking.

\subsection{Optimization ablation}

The optimization campaign reduced the best observed subtree-parallel runtime
approximately as follows:
\[
1.2530
\rightarrow
1.2007
\rightarrow
1.1619
\rightarrow
1.0179
\rightarrow
0.93533
\rightarrow
0.89257
\quad\mathrm{ms}.
\]

The final value is approximately
\[
28.8\%
\]
below the earlier \(1.2530\)-ms subtree baseline.

The major experimentally validated improvements were:

\begin{enumerate}
    \item specialized Goldilocks multiplication reduction;
    \item fused field-level kernel fold;
    \item elimination of bounds checks in the hot level loop;
    \item direct first-level reduction into half-size scratch;
    \item worker-local scratch reuse.
\end{enumerate}

The isolated microbenchmarks showed approximately:

\begin{center}
\begin{tabular}{lrrr}
\toprule
Optimization &
Before &
After &
Local improvement \\
\midrule
Goldilocks multiplication &
\(2.2372\) ns &
\(2.0559\) ns &
\(8.1\%\) \\

Fused kernel fold &
\(4.0591\) ns &
\(3.3585\) ns &
\(17.3\%\) \\

Unchecked level loop &
\(6.9671\,\mu\mathrm{s}\) &
\(6.3575\,\mu\mathrm{s}\) &
\(8.75\%\) \\

Direct first level &
\(6.9031\,\mu\mathrm{s}\) &
\(6.1571\,\mu\mathrm{s}\) &
\(10.8\%\) \\
\bottomrule
\end{tabular}
\end{center}

These percentages describe the affected primitive or stage and should not be
added together to predict whole-evaluator speedup.

\subsection{Negative results}

Two-level and three-level tree fusion were both rejected.  In particular,
the two-level experiment gave
\[
9.7019\ \mu\mathrm{s}
\]
for the ordinary two-level path and
\[
9.8528\ \mu\mathrm{s}
\]
for the fused path.

Thus the version with fewer intermediate memory accesses was approximately
\(1.6\%\) slower.

This result illustrates an important implementation constraint: reducing
memory traffic does not necessarily improve scalar field code when doing so
increases dependency depth and register pressure.

\subsection{Interpretation}

The final results establish a narrow but concrete claim:

\begin{quote}
For arbitrary single-point evaluation of a degree-\(<N\) univariate
polynomial already stored in the tested representation, the optimized
Boolean-kernel evaluator is faster than the monomial and roots-of-unity
Lagrange implementations in our Rust benchmark over Goldilocks on the
tested Apple M1 machine.
\end{quote}

At \(N=2^{20}\), the conservative same-run comparison gives a
\[
1.337\times
\]
advantage over the parallel monomial baseline and a
\[
25.2\%
\]
reduction in wall-clock evaluation time.

The experiment does not yet establish an end-to-end prover speedup.  That
requires measuring basis conversion, encoding, commitment, FRI-related
operations, hashing, and all other protocol costs in an integrated system.
